% Evidence-Certificate GraphRAG: method section for the Graph-R1 paper.
% Compile standalone:  latexmk -pdf evidence_certificate.tex
% To include in the paper, copy everything between \begin{document} and \end{document} (or \input it).
% Every definition below is the one implemented in evidence/*.py (file names are given in the text).
\documentclass[10pt]{article}
\usepackage[margin=2.2cm]{geometry}
\usepackage{amsmath,amssymb,booktabs,algorithm,algpseudocode,enumitem,xcolor}
\newcommand{\Risk}{\mathrm{Risk}}
\newcommand{\VNE}{\mathrm{VNE}}
\newcommand{\GE}{G_E}
\newcommand{\GH}{G_H}
\newcommand{\todo}[1]{\textcolor{red}{[TODO: #1]}}
\begin{document}

\section{Evidence-Certificate Reasoning on the Knowledge Hypergraph}
\label{sec:evidence}

Graph-R1 trains an agent with reinforcement learning to query the knowledge hypergraph $\GH$ and to
decide when it knows enough. Its stopping decision is made by the policy itself, so an answer carries
no statement of \emph{what evidence supports it} and the agent cannot tell a lucky guess from a fact that is
documented. We add a \emph{training-free} layer on top of the unchanged $\GH$ and the unchanged dual-path
retrieval API. A frozen LLM (no gradient update) is used only to (i) decompose the question and (ii) match
passages to facts. Everything that decides \emph{which evidence to buy next} and \emph{when to stop} is an explicit,
inspectable procedure over an \emph{evidence support graph}, and the output is an \emph{evidence certificate}.

\paragraph{Overview.} For a question $q$ the system runs the loop
\[
q \;\to\; \mathcal O(q) \;\to\; \text{evidence state } s_t \;\to\; \text{audit}(s_t) \;\to\; \text{controller} \;\to\; \text{retrieval on }\GH \;\to\; s_{t+1}\;\ldots\;\to\; (a,\text{cert}).
\]

\subsection{Evidence obligations}
The LLM maps $q$ to a set $\mathcal O(q)=\{o_1,\dots,o_n\}$, $n\le 6$ (one call, temperature $0$, strict JSON, up to three
validated attempts; \texttt{obligations.py}). An obligation $o_i=(\tau_i,\ \mathrm{dep}_i,\ \kappa_i)$ is an atomic statement $\tau_i$ with
unknown variables $?X$, a dependency list $\mathrm{dep}_i\subseteq\mathcal O$ and a kind
$\kappa_i\in\{\textsf{retrieve},\textsf{derive}\}$. A \textsf{retrieve} obligation has to be stated by a document of the
corpus; a \textsf{derive} obligation (compare two dates, same country?) is computed from the values of others and needs no retrieval.
One designated obligation $o_{\mathrm{ans}}$ produces the answer. The dependencies form a DAG (cycles are rejected). Validation also
requires: every variable is introduced by exactly one obligation (otherwise values of unrelated unknowns would merge); every
parenthesised qualifier of the question and the type word of a named entity (\emph{film}, \emph{album}, \dots) survive in the obligations; and
$o_{\mathrm{ans}}$ depends, directly or not, on at least one \textsf{retrieve} obligation. Let $R\subseteq\mathcal O$ be the required
\textsf{retrieve} obligations (the ancestors of $o_{\mathrm{ans}}$ of kind \textsf{retrieve}).

\subsection{Evidence state and provenance}
The state $s$ holds the retrieved passages (items), for each obligation the set of items that \emph{support} it with the variable
bindings the matcher extracted, and the values of the variables inherited at match time. A match is accepted only if all of these hold
(\texttt{state.validate\_match}): the quoted words occur in the passage (at least 10 normalised characters); every bound value occurs in the passage;
the value has the type its slot asks for (a date for \emph{born on}, not a date for \emph{born in} when the question asks \emph{where});
the value is not the entity the question is about nor equal to the value it was derived from (``the father of $X$ is $X$''); and the passage is
about the entity the chain arrived at (subject check). The matcher also receives the question, so that passages about a same-named entity of
another kind (an album and a film) are rejected.

\paragraph{Provenance unit.} The independence of two sources is judged on the \emph{article} they come from, not on the passage
(\texttt{provenance.py}): each passage is mapped back, through the chunk store of the graph, to the title(s) of the article it was extracted from.
Passages that share a title are \emph{one} unit (union--find over titles). A passage whose article is unknown is not an independent source: it counts
only if nothing with a known title supports the fact. For a variable $?X$ the \emph{chosen value} is the supported value with the most
independent units (ties by retrieval coherence); alias forms of one name (``Annakin'', ``Ken Annakin'') are one value, regnal numbers, place suffixes
(``of Wales''/``of Prussia'') and parenthesised qualifiers keep entities apart. A support entry is \emph{live} if it was matched under the
\emph{current} values of the inherited variables, agrees with the chosen value of the obligation's own variables, and its unit is not removed.

\subsection{Evidence support graph and audit}
The evidence support graph $\GE=(V,E)$ is a typed, directed graph with nodes for required obligations, live items and provenance units,
and edges \emph{depends} ($o_j\!\to\! o_i$ for $o_j\in\mathrm{dep}_i$), \emph{supports} (item $\to$ obligation) and \emph{from} (unit $\to$ item).
For $o\in R$ let $U(o)$ be the set of independent units connected to $o$ and $\mathrm{k}(o)=|U(o)|$ ($=0$ if uncovered).
\begin{itemize}[leftmargin=1.2em,itemsep=1pt]
 \item \textbf{Coverage} $C=\frac1{|R|}\,|\{o\in R:\mathrm{k}(o)\ge1\}|$.
 \item \textbf{Critical cuts.} A unit $u$ is a cut if $U(o)=\{u\}$ for some $o\in R$ (its removal disconnects $o$ from all evidence). A cut is
  \emph{audited} once the Verify action confirmed that its sole passage states the fact (see below).
 \item \textbf{Robustness} $R_{\mathrm{rob}}=\frac{|\{o \text{ covered}: \mathrm{k}(o)\ge2\ \vee\ o\text{ audited}\}|}{|\{o\text{ covered}\}|}$.
 \item \textbf{Provenance diversity} $D=\frac1{|R|}|\{o\in R: \text{at least two known units}\}|$.
 \item \textbf{Conflicts.} A variable with several supported values is in conflict; it is \emph{resolved} iff the best value has strictly more known
  independent units than the runner-up. $X$ is the number of unresolved conflicts.
 \item \textbf{Counterfactual fragility.} For every unit $u$ the state is recomputed without $u$; $u$ is fragile if some obligation becomes uncovered
  or the chosen value of a variable changes. The list is reported in the certificate and drives the Diversify gain.
\end{itemize}
The \emph{state vector} is $\mathbf v(s)=(C,\ R_{\mathrm{rob}},\ D,\ -X)$; a step \emph{improves} the state if no component decreases and one increases (Pareto
test). A scalar alternative is $\Risk(s)=(1-C)+\lambda_X X+\lambda_R(1-R_{\mathrm{rob}})$ ($\lambda_X=\lambda_R=0.5$). A step that \emph{discovers} a conflicting value
(Diversify, Challenge) is counted as a success although $X$ grows, because that evidence is what the audit needs.

\paragraph{Certificate.} $\mathrm{Cert}(s)$ is \emph{complete} iff (i) $R\neq\emptyset$ and every $o\in R$ is covered, (ii) the support graph over the required
obligations is connected, (iii) no conflict is unresolved, (iv) every critical cut is audited. A complete certificate is attached to the answer only if, in
addition, the \emph{answer is grounded}: for a \textsf{retrieve} answer the answer string must be (an alias of) the chosen value of the answer
obligation's own variable, or occur verbatim in a passage supporting it (a value that belongs to another obligation does not count); for a \textsf{derive}
answer it must be \emph{computed} (below) or one of the options named in the question, or yes/no. The certificate records the basis (\emph{retrieved},
\emph{computed}, \emph{option}, \emph{yes/no}); only the first two check the final step against the evidence, the last two certify the facts and leave the
comparison to the LLM. Model confidence plays no role anywhere.

\paragraph{Deterministic comparison.} For ``which of $A$ or $B$ \ldots earlier/later'' questions with two date variables, the last step is computed, not
generated (\texttt{derive.py}): each date is traced back along the obligation chain to the option it belongs to, dates are compared at the coarsest
shared precision, and the function returns nothing (the LLM answer is kept) if the options, the direction or a tie are ambiguous.

\subsection{Decision-time controller}
Each step the controller enumerates the applicable actions, prices them, and executes the one with the best \emph{value of next evidence}
\[
\VNE(a\mid s)\;=\;p_a\cdot g_a\;-\;\mathrm{cost}(a),\qquad \mathrm{cost}(a)=c_s\,n_{\mathrm{search}}(a)+c_\ell\,n_{\mathrm{llm}}(a).
\]
\begin{table}[h]\centering\small
\begin{tabular}{@{}p{1.5cm}p{3.7cm}p{6.6cm}p{3.1cm}@{}}\toprule
action & applicable when & what it does & gain $g_a$\\ \midrule
FillGap & required $o$ uncovered and ready & search its fact (plain query, compact form, entity form, LLM paraphrase, then up to 2 \emph{informed} LLM queries that read the retrieved titles) & $w_{gap}(1+\#\text{blocked descendants})$\\
Bridge & upstream hop of $o$ cannot be found & search the chain as one query and the whole question & $w_{gap}(1+\cdot)$\\
Diversify & $o$ covered with $<2$ known units, or tied conflict & search a second article stating the fact (one query) & $w_{fragile}$ / $w_{conflict}$\\
Challenge & conflict, or confirmation bias & search for passages contradicting the chosen value & $w_{conflict}$ / $w_{chal}$\\
Verify & critical cut not audited & re-check the sole passage with a stricter LLM prompt; remove the support if it fails (no search) & $w_{unver}$\\
Answer & no action with $\VNE>0$ & stop & --\\ \bottomrule
\end{tabular}
\end{table}
$p_a$ is the estimated success probability of the action kind, $p_a=(s_a+\alpha_0)/(n_a+\alpha_0+\beta_0)$ with a Beta prior ($\alpha_0=1,\beta_0=3$, $p_0=0.25$) and
success meaning a Pareto improvement. The counts $(s_a,n_a)$ are carried over the questions of a run, so an optional action that does not pay on the corpus
is dropped automatically; FillGap and Verify never fall below the prior (a certificate needs them) and a conflict-driven Diversify keeps its own counts.
The gains only fix an order of importance (gap $>$ conflict $>$ challenge $>$ fragile $>$ unverified) and are \emph{not} an exact expected risk reduction; the
Pareto test decides success. Defaults: $w_{gap}=1,\ w_{conflict}=0.8,\ w_{chal}=0.2,\ w_{fragile}=0.3,\ w_{unver}=0.15,\ c_s=0.05,\ c_\ell=0.02$.
For every executed step we also log the shaped reward $r=\Delta C+\alpha\Delta R_{\mathrm{rob}}+\beta\Delta D-\gamma\Delta X-\lambda\,\mathrm{cost}$
($\alpha=\beta=\gamma=0.5,\ \lambda=1$). It is a diagnostic: the controller is training-free and no policy is optimised on it.

\paragraph{Reasoning-pathology monitor.} (i) \emph{retrieval loop}: three consecutive search steps without a new passage or a Pareto gain stop the search;
(ii) \emph{over-retrieval}: two searches for an obligation without a new passage close it; (iii) \emph{confirmation bias}: every covered obligation rests on one known
article and nothing was challenged, so Challenge is allowed for the weakest one; (iv) \emph{premature confidence}: an answer produced while a required obligation is uncovered;
(v) \emph{ungrounded answer}: the answer occurs in none of the evidence; (vi) \emph{derive mismatch}: the LLM's comparison disagrees with the computed one.
(i)--(iii) act on the control flow, (iv)--(vi) are reported.

\paragraph{Stopping.} The loop ends when the certificate is complete and no action has $\VNE>0$ (\emph{certified}), when no action has $\VNE>0$ but the
certificate is incomplete, when upstream hops are blocked, on a retrieval loop, or when the search budget (10 queries) or the step budget (12) is exhausted.
An uncertified question is still answered (flagged) so that EM can be reported both ways; \texttt{--abstain-uncertified} refuses instead.

\begin{algorithm}[h]\caption{Evidence-certificate controller for one question}\small
\begin{algorithmic}[1]
\State $\mathcal O\gets\textsc{Decompose}(q)$; $s\gets\textsc{InitState}(\mathcal O)$
\For{$t=1\ldots T_{\max}$}
 \State \textsc{Sweep}: show held, unchecked passages to the matcher for obligations that just became ready
 \State $\mathrm{cert}\gets\textsc{Audit}(s)$; $\mathrm{flags}\gets\textsc{Monitor}(s)$
 \If{flags.retrieval\_loop} \textbf{break} \EndIf
 \State $\mathcal A\gets\textsc{Propose}(s,\mathrm{cert},\mathrm{flags})$; $a^\star\gets\arg\max_{a\in\mathcal A}\VNE(a)$
 \If{$\mathcal A=\emptyset$ \textbf{or} $\VNE(a^\star)\le0$ \textbf{or} budget exceeded} \textbf{break} \EndIf
 \State execute $a^\star$ (search $\GH$ / verify), match new passages, update $s$; update $p_{a^\star}$ with the Pareto test
\EndFor
\State $a\gets\textsc{Answer}(s)$ (computed for date comparisons, otherwise from the established facts)
\State \Return $a$, $\mathrm{Cert}(s)$, certified $\;=\;\mathrm{Cert}$ complete $\wedge$ answer grounded
\end{algorithmic}\end{algorithm}

\subsection{Implementation details and cost}
The hypergraph, the embeddings and the dual-path retrieval (entity and hyperedge, rank fusion) are used through the existing \texttt{/search} API
without change. One matcher call sees at most 10 passages of at most 700 characters. The audit is polynomial (graph of $O(n+|\text{items}|)$ nodes; fragility
recomputes the state once per unit). Per question the controller issues at most 10 retrieval queries and, in our runs, about 10 LLM calls (decomposition, matching, verification,
query writing, answer). All thresholds above are hyper-parameters on the command line.

\subsection{Assumptions and limits}
(1) A fact of the corpus is traceable to a source article; the provenance index is built from the graph's chunk store and an item without a title is never counted as independent
(the controller warns when most supporting passages have no title). (2) Type cues, date parsing and name aliases are English; the matcher and obligation prompts are not
dataset-specific but were tuned on 2WikiMultiHopQA. (3) The certificate covers retrieved facts; it cannot detect a passage about a different entity with the same name
unless the question or the chain disambiguates it, and the final comparison is computed only for dates. (4) Independent second sources are rare in 2WikiMultiHopQA, so $D$ stays low there. (5) The value
of next evidence is an estimate with hand-set gains, not an exact expected risk reduction.

\subsection{Experimental protocol}
\textbf{Systems} (same frozen LLM, same retrieval API, same questions): the Graph-R1 agent loop without RL (baseline), the controller, and ablations: without Diversify, Verify, Bridge, Challenge;
FillGap only; without the fragility term; guards off; every passage independent (naive provenance); scalar risk; LLM comparison; per-question $p$.
\textbf{Metrics}: EM, F1, retrieval queries and LLM calls per question; the share of certified answers; EM on certified and uncertified answers; selective EM
(answer only when certified, abstention counts as wrong). \textbf{Statistics}: Wilson 95\% intervals; paired wins/losses against a reference, exact McNemar and a bootstrap interval of the EM difference
(\texttt{evidence/compare.py}). \textbf{Data}: a development set on which the guards were tuned and a disjoint held-out set (\texttt{--exclude}).
\todo{insert final table from runs/final/report.md}

\end{document}
